\uuid{JViK}
\titre{Logique}
\auteur{bodin ; croizat ; sacre}
\organisation{exo7}
\licence{CC BY-NC-SA 4.0}
\datecreate{2022-03-08}
\niveau{L1}
\module{Algèbre}
\chapitre{Logique, ensemble, raisonnement}
\sousChapitre{Logique}
\difficulte{3}
\theme{Logique}
\exo7qcmid{hilisit.logique.015}

\contenu{
\begin{question}
\qclassification{100.01}
Soit $\mathcal{P}$ et $\mathcal{Q}$ deux assertions telles que "$\mathcal{P} \implies \mathcal{Q}$" soit vraie, et "$\text{non}(\mathcal{P}) \implies \text{non}(\mathcal{Q})$" soit aussi vraie. On a alors :
\begin{answers}
	\good{"$\mathcal{Q} \implies \mathcal{P}$" est vraie.}
	\good{"$\mathcal{P} \iff \mathcal{Q}$" est vraie.}
	\bad{"$\mathcal{Q} \implies \text{non}(\mathcal{P})$" est vraie.}
	\bad{"$\mathcal{P} \implies \text{non}(\mathcal{Q})$" est vraie.}
\end{answers}
\begin{explanations}
Par contraposition de "$\text{non}(\mathcal{P}) \implies \text{non}(\mathcal{Q})$", on obtient "$\mathcal{Q} \implies \mathcal{P}$". On a donc "$\mathcal{P} \iff \mathcal{Q}$". Ceci signifie que les assertions $\mathcal{P}$ et $\mathcal{Q}$ sont soit simultanément vraies, soit simultanément fausses. Cela exclut la possibilité d'avoir l'une vraie et l'autre fausse.
\end{explanations}
\end{question}
}
